\documentclass[10pt,a4paper]{amsart}
\usepackage[margin=19mm]{geometry}
\usepackage{amsmath,amssymb,mathtools}
\usepackage[scr=rsfs,cal=euler]{mathalfa}
\usepackage{xfrac}
\usepackage[unicode,hidelinks,bookmarks=false]{hyperref}
\newcommand{\ie}{i.e.\ }
\newcommand{\cf}{cf.\ }
\newcommand{\resp}{respectively\ }
\newcommand{\Subsection}{\subsection}
\providecommand{\nobreakdash}{\nobreak\hbox{-}}
\setlength{\emergencystretch}{2em}
\multlinegap0pt
\usepackage{mathrsfs}
\usepackage{enumitem}
\newtheorem{theorem}{Theorem}[section]
\newtheorem{lemma}[theorem]{Lemma}
\newtheorem{proposition}[theorem]{Proposition}
\newtheorem{corollary}[theorem]{Corollary}
\newtheorem{keyword}[theorem]{Keywords}
\newtheorem{definition}[theorem]{Definition}
\newtheorem{remark}[theorem]{Remark}
\newtheorem{example}[theorem]{Example}
\begin{document}
\title[On Locally Dually and Projectively Flat Almost Rational Fins]{On Locally Dually and Projectively Flat Almost Rational Finsler Metrics}
\author{Ghanashyam Kumar Prajapati, Ranadip Gangopadhyay, Bankteshwar Tiwari}
\date{}
\maketitle
\noindent\emph{Source and licence.} On Locally Dually and Projectively Flat Almost Rational Finsler Metrics. Version arXiv:2606.15699v1 (CC BY 4.0). This material is distributed under the Creative Commons Attribution 4.0 International licence, http://creativecommons.org/licenses/by/4.0/, as recorded in the arXiv OAI licence metadata for this exact version and shown on the abstract page https://arxiv.org/abs/2606.15699v1. Full rights notice: licenses/arxiv-2606.15699-CC-BY-4.0.txt.\par\smallskip
\noindent\textit{Editorial scope.} Counted unit: the original \texttt{theorem} environment \texttt{4.2} at source lines 510--524 together with its complete proof at lines 526--630; the delivered document keeps source lines 479--631 so that every referenced label and every citation resolves inside it. Native editable source is the paper's own concatenated TeX; the AMS wrapper is editorial formatting only. No publisher class, upstream script or unrelated artwork is redistributed.
								\section{Locally Projectively Flat AR-Finsler Metrics} 
								
								Projectively flat Finsler metrics are characterized by having straight-line geodesics in suitable local coordinates. While in Riemannian geometry projective flatness is equivalent to constant sectional curvature, Finsler geometry admits genuinely non-Riemannian examples such as the $m$-th root Finsler metrics. These rigidity properties motivate the study of projective flatness in the broader class of Almost Rational (AR) Finsler metrics, where rationality is imposed only on selected geometric objects instead of the whole fundamental tensor.\\
								It is known \cite{hamel1903} that a Finsler metric is projectively flat if and only if
								\begin{equation}\label{2.3e}
								F_{x^k y^l} y^k - F_{x^l} = 0.
								\end{equation}
								
								\begin{lemma}\label{lem1}
									If $g_{ij}(x,y)$ is the fundamental tensor of a Finsler metric, then
									\[
									(g_{ij})_{x^k y^l}y^i y^j y^k
									=
									-(g_{lj})_{x^k}y^j y^k.
									\]
								\end{lemma}
								
								\begin{proof}
									Since $g_{ij}$ is $0$-homogeneous in $y$,
									$(g_{ij})_{y^l}y^l=0.$
									Differentiating with respect to $x^k$ gives
									\[
									(g_{ij})_{x^k y^l}y^l=-(g_{ij})_{x^k}.
									\]
									Multiplying by $y^i y^j$ and using the symmetry of $g_{ij}$ together with a relabelling of dummy indices yields
									\[
									(g_{ij})_{x^k y^l}y^i y^j y^k
									=
									-(g_{lj})_{x^k}y^j y^k.
									\]
								\end{proof}

								\begin{theorem}\label{4.2} Let $F$ be  an almost rational Finsler metric, then $F$ is
									projectively flat 
									if it satisfies the following condition:
									
									\begin{equation}\label{eq:PF_AR_reduced}
									-\frac{F_{y^l}}{F}
									(\eta a_{ij})_{x^k}y^iy^jy^k
									+
									(\eta a_{lj})_{x^k}y^jy^k
									-
									(\eta a_{ij})_{x^l}y^iy^j
									=0.
									\end{equation}
									
								\end{theorem}

								\begin{proof}
									Since
									\[
									F^2=\eta a_{ij}y^iy^j,
									\]
									differentiation with respect to \(x^l\) yields
									\[
									2FF_{x^l}
									=
									(\eta a_{ij})_{x^l}y^iy^j.
									\]
									Hence
									\begin{equation}
									F_{x^l}
									=
									\frac{1}{2F}
									(\eta a_{ij})_{x^l}y^iy^j.
									\label{eq:Fxl}
									\end{equation}
									
									Differentiating \eqref{eq:Fxl} with respect to \(y^l\), we obtain
									\[
									F_{x^ky^l}
									=
									-\frac{F_{y^l}}{2F^2}
									(\eta a_{ij})_{x^k}y^iy^j
									+\frac{1}{2F}
									\frac{\partial}{\partial y^l}
									\Big(
									(\eta a_{ij})_{x^k}y^iy^j
									\Big).
									\]
									
									Since
									\[
									\frac{\partial}{\partial y^l}
									\Big(
									(\eta a_{ij})_{x^k}y^iy^j
									\Big)
									=
									(\eta a_{ij})_{x^ky^l}y^iy^j
									+
									2(\eta a_{lj})_{x^k}y^j,
									\]
									it follows that
									\[
									F_{x^ky^l}
									=
									\frac{1}{2F}
									\left[
									-\frac{F_{y^l}}{F}
									(\eta a_{ij})_{x^k}y^iy^j
									+
									(\eta a_{ij})_{x^ky^l}y^iy^j
									+
									2(\eta a_{lj})_{x^k}y^j
									\right].
									\]
									
									Multiplying by \(y^k\), we get
									\[
									F_{x^ky^l}y^k
									=
									\frac{1}{2F}
									\left[
									-\frac{F_{y^l}}{F}
									(\eta a_{ij})_{x^k}y^iy^jy^k
									+
									(\eta a_{ij})_{x^ky^l}y^iy^jy^k
									+
									2(\eta a_{lj})_{x^k}y^jy^k
									\right].
									\]
									
									Since \(F\) is projectively flat, \eqref{2.3e}
									holds. Substituting the above expressions and multiplying by \(2F\), we obtain
									\begin{equation}\label{26}
									-\frac{F_{y^l}}{F}
									(\eta a_{ij})_{x^k}y^iy^jy^k
									+
									(\eta a_{ij})_{x^ky^l}y^iy^jy^k
									+
									2(\eta a_{lj})_{x^k}y^jy^k
									-
									(\eta a_{ij})_{x^l}y^iy^j
									=0.   
									\end{equation}
									
									Now from Lemma \ref{lem1}, we have
									$(\eta a_{ij})_{x^ky^l}y^iy^jy^k
									=
									-(\eta a_{lj})_{x^k}y^jy^k.$
									Substituting this identity into \eqref{26} gives
									\[
									-\frac{F_{y^l}}{F}
									(\eta a_{ij})_{x^k}y^iy^jy^k
									+
									(\eta a_{lj})_{x^k}y^jy^k
									-
									(\eta a_{ij})_{x^l}y^iy^j
									=0.
									\]
									
									Hence \eqref{eq:PF_AR_reduced} follows.
								\end{proof}
								

\begin{thebibliography}{99}
\bibitem[hamel1903]{hamel1903} G. Hamel, Über die Geometrien, in denen die Geraden die Kürzesten sind, Math. Ann., 57 (1903), 231--264.
\end{thebibliography}
\end{document}
